%========================= BIBLIOGRAPHIE / WEBOGRAPHIE ======================
% Chapitre non numéroté : une seule entrée dans la table des matières (via
% \addcontentsline) et marques d'en-tête posées explicitement. La liste est
% composée manuellement (pas d'environnement thebibliography) afin d'éviter
% le double titre « Bibliographie » et l'en-tête erroné sur les pages
% suivantes que cet environnement génère dans la classe report.
\chapter*{Bibliographie et webographie}
\addcontentsline{toc}{chapter}{Bibliographie et webographie}
\markboth{Bibliographie et webographie}{Bibliographie et webographie}

Les références ci-dessous renvoient aux documentations officielles des
technologies utilisées dans le projet ; il s'agit d'une webographie, aucune
source académique n'étant nécessaire au sujet traité. L'ensemble de ces
ressources en ligne a été consulté le 11~juillet~2026. Une version
\texttt{BibTeX} de ces références est fournie dans le fichier
\texttt{references.bib} accompagnant ce rapport.

\begin{enumerate}[label={[\arabic*]}, leftmargin=1.5cm, itemsep=4pt]
  \item React --- \emph{Documentation officielle}.
        \url{https://react.dev/}
  \item Vite --- \emph{Next Generation Frontend Tooling}.
        \url{https://vitejs.dev/}
  \item React Router --- \emph{Documentation officielle}.
        \url{https://reactrouter.com/}
  \item FastAPI --- \emph{Documentation officielle}.
        \url{https://fastapi.tiangolo.com/}
  \item Pydantic --- \emph{Data validation using Python type hints}.
        \url{https://docs.pydantic.dev/}
  \item OpenAPI Initiative --- \emph{OpenAPI Specification}.
        \url{https://www.openapis.org/}
  \item Swagger --- \emph{API Documentation \& Design Tools}.
        \url{https://swagger.io/}
  \item SQLAlchemy --- \emph{The Database Toolkit for Python}.
        \url{https://www.sqlalchemy.org/}
  \item PostgreSQL --- \emph{The World's Most Advanced Open Source Relational
        Database}. \url{https://www.postgresql.org/}
  \item Leaflet --- \emph{An open-source JavaScript library for interactive
        maps}. \url{https://leafletjs.com/}
  \item OpenStreetMap --- \emph{Carte collaborative libre}.
        \url{https://www.openstreetmap.org/}
  \item Recharts --- \emph{A composable charting library built on React
        components}. \url{https://recharts.org/}
  \item Tailwind CSS --- \emph{A utility-first CSS framework}.
        \url{https://tailwindcss.com/}
  \item GitHub --- \emph{Plateforme d'hébergement et de collaboration Git}.
        \url{https://github.com/}
\end{enumerate}
